\documentclass[tikz,border=8pt]{standalone}
\usepackage{tikz}
\usepackage{amsmath}
\usepackage{amssymb}
\usepackage{pgfplots}
\pgfplotsset{compat=1.18}
\usepackage{ctex}
\usetikzlibrary{calc,positioning,arrows.meta,matrix}

% LC 1143 最长公共子序列（LCS）
% s1 = "abcde"（行，i=0..5）
% s2 = "ace"  （列，j=0..3）
% dp[i][j] = LCS(s1[0..i-1], s2[0..j-1])
% 答案 = dp[5][3] = 3 ("ace")

\begin{document}
\begin{tikzpicture}[
  every node/.style={font=\small},
  cell/.style={draw=gray!70, minimum width=0.9cm, minimum height=0.78cm,
               inner sep=0pt, thick, fill=white},
  initcell/.style={draw=blue!50, minimum width=0.9cm, minimum height=0.78cm,
                   inner sep=0pt, thick, fill=blue!8},
  tracecell/.style={draw=red!60, minimum width=0.9cm, minimum height=0.78cm,
                    inner sep=0pt, very thick, fill=red!6},
  rescell/.style={draw=orange!70, minimum width=0.9cm, minimum height=0.78cm,
                  inner sep=0pt, very thick, fill=orange!8},
  arr/.style={->,>=Stealth,thick},
  tracearr/.style={->,>={Stealth[scale=0.6]},thick,red!70},
]

%% ── Layout constants ──────────────────────────────────────
% cs=1.1, rs=0.92, ox=0.5, oy=-0.8
% cell center: x = ox + j*cs + 0.45,  y = oy - i*rs - 0.39
% j=3 right edge = 0.5 + 3*1.1 + 0.45 + 0.45 = 4.70
% Table y: top=-1.19, bottom=-6.19

\def\cs{1.1}
\def\rs{0.92}
\def\ox{0.5}
\def\oy{-0.8}

%% ── 标题 ──────────────────────────────────────────────────
% col headers now at oy+0.65=-0.15; title at +0.55 and +1.2
\node[font=\large\bfseries] at (2.2, 1.3)
  {最长公共子序列 LCS（LC 1143）};
\node[font=\small, gray] at (2.2, 0.7)
  {s1="abcde"，s2="ace"，dp[i][j] = LCS 长度，红色 = 回溯路径};

%% ── 列标签（s2: ∅ a c e）
% oy=-0.8; table top at oy. Characters at oy+0.65, j-index at oy+0.38
% => j-index is 0.38cm above table top — clear of all cells
\foreach \j/\ch in {0/$\varnothing$, 1/a, 2/c, 3/e} {
  \node[font=\small\bfseries] at (\ox+\j*\cs+0.45, \oy+0.65) {\ch};
  \node[font=\tiny, gray]     at (\ox+\j*\cs+0.45, \oy+0.38) {j=\j};
}

%% ── 行标签（s1: ∅ a b c d e）
\foreach \i/\ch in {0/$\varnothing$, 1/a, 2/b, 3/c, 4/d, 5/e} {
  \node[font=\small\bfseries, anchor=east] at (\ox-0.12, \oy-\i*\rs-0.39) {\ch};
  \node[font=\tiny, gray, anchor=east]    at (\ox-0.50, \oy-\i*\rs-0.39) {i=\i};
}

%% ── dp 表格数值 ────────────────────────────────────────────
\foreach \i/\j/\v/\style in {
  % row 0 (init)
  0/0/0/initcell, 0/1/0/initcell, 0/2/0/initcell, 0/3/0/initcell,
  % col 0 (init)
  1/0/0/initcell, 2/0/0/initcell, 3/0/0/initcell, 4/0/0/initcell, 5/0/0/initcell,
  % trace cells (match positions in red)
  1/1/1/tracecell,
  3/2/2/tracecell,
  5/3/3/rescell,
  % non-trace interior
  1/2/1/cell, 1/3/1/cell,
  2/1/1/cell, 2/2/1/cell, 2/3/1/cell,
  3/1/1/cell,              3/3/2/cell,
  4/1/1/cell, 4/2/2/cell, 4/3/2/cell,
  5/1/1/cell, 5/2/2/cell%
} {
  \node[\style] (D\i\j)
    at (\ox+\j*\cs+0.45, \oy-\i*\rs-0.39) {\v};
}

%% ── 回溯路径箭头（红色虚线短箭头，从右下往左上）
%% 仅在相邻 cell 间走 corner→corner，使用 shorten 让箭头不触及 cell 边框
\draw[tracearr, dashed, shorten <=2pt, shorten >=2pt]
  (D53.north west) -- (D42.south east);  % (5,3)→(4,2) e=e
\draw[tracearr, dashed, shorten <=2pt, shorten >=2pt]
  (D42.north)      -- (D32.south);       % (4,2)→(3,2) d≠c
\draw[tracearr, dashed, shorten <=2pt, shorten >=2pt]
  (D32.north west) -- (D21.south east);  % (3,2)→(2,1) c=c
\draw[tracearr, dashed, shorten <=2pt, shorten >=2pt]
  (D21.north)      -- (D11.south);       % (2,1)→(1,1) b≠a
\draw[tracearr, dashed, shorten <=2pt, shorten >=2pt]
  (D11.north west) -- (D00.south east);  % (1,1)→(0,0) a=a

%% ── 匹配点标注移除：表内仅保留红色 cell + 虚线箭头标识路径
%% 完整的路径序列 a@(1,1), c@(3,2), e@(5,3) 已在底部"完整 dp 表"框注明

%% ── 分隔线（表格底部 y≈-6.24，间距 0.5cm）
\draw[gray!30, dashed] (-0.6, -6.7) -- (7.0, -6.7);

%% ── 转移公式框（居中于表格，表格下方）
\node[draw=gray!50, fill=white, rounded corners=3pt, font=\small,
      inner sep=8pt, align=left]
  at (2.4, -7.9)
  {\textbf{转移公式}\quad
   若 \texttt{s1[i]}$=$\texttt{s2[j]}：
   $\;\text{dp}[i][j]=\text{dp}[i\!-\!1][j\!-\!1]+1$\\[5pt]
   否则：$\;\text{dp}[i][j]=\max\!\bigl(\text{dp}[i\!-\!1][j],\;\text{dp}[i][j\!-\!1]\bigr)$\\[5pt]
   \textbf{回溯：}从 dp[5][3] 沿红色箭头追踪对角匹配，还原 LCS = \textbf{"ace"}};

%% ── 结果框
\node[draw=gray!50, fill=white, rounded corners=3pt, font=\small,
      inner sep=7pt, align=left]
  at (2.4, -10.5)
  {\textbf{完整 dp 表：}\\[3pt]
   \quad$j{=}0$: [0,0,0,0,0,0]\quad
   $j{=}1$: [0,1,1,1,1,1]\quad
   $j{=}2$: [0,1,1,2,2,2]\quad
   $j{=}3$: [0,1,1,2,2,3]\\[4pt]
   \textbf{LCS = 3}，公共子序列为 \textbf{"ace"}
   \quad（红色路径：a$@$(1,1)，c$@$(3,2)，e$@$(5,3)）\\
   \textbf{复杂度：}时间 $O(m{\times}n)$，空间 $O(m{\times}n)$ 或 $O(n)$ 滚动数组};

%% ── 图例
\node[draw=gray!50, fill=white, rounded corners=3pt, font=\scriptsize,
      inner sep=4pt, align=left]
  at (2.4, -12.5)
  {\textbf{图例：}
   蓝色 = 初始化（空串）\quad
   红色格+箭头 = 回溯匹配路径\quad
   橙色 = 最终答案格};

\end{tikzpicture}
\end{document}
